\documentclass{beamer}
\usepackage{graphicx}
\usepackage{amsmath, amssymb}
\usepackage{tikz}
\usetikzlibrary{positioning}
\usepackage{array}

\title{Assignment}
\author{Akhil Mohan}
\date{\today}

\begin{document}

% Title Slide
\begin{frame}
\titlepage
\end{frame}

% Examples of Sets
\begin{frame}{Examples of Sets}
Even numbers:


\[
A = \{ x \in \mathbb{Z} \mid x \text{ is even} \}
\]



Odd numbers:


\[
B = \{ x \in \mathbb{Z} \mid x \text{ is odd} \}
\]


\end{frame}

\begin{frame}{Examples of Sets (contd.)}
Natural numbers less than 10:


\[
C = \{ x \in \mathbb{N} \mid x < 10 \}
\]



Multiples of 5:


\[
D = \{ x \in \mathbb{Z} \mid x = 5n, \, n \in \mathbb{Z} \}
\]


\end{frame}

\begin{frame}{Examples of Sets (contd.)}
Prime numbers:


\[
E = \{ x \in \mathbb{N} \mid x \text{ is prime} \}
\]



Integers between -5 and 5:


\[
F = \{ x \in \mathbb{Z} \mid -5 \leq x \leq 5 \}
\]


\end{frame}

\begin{frame}{Examples of Sets (contd.)}
Perfect squares:


\[
G = \{ x \in \mathbb{N} \mid x = n^2, \, n \in \mathbb{N} \}
\]



Real numbers greater than 0:


\[
H = \{ x \in \mathbb{R} \mid x > 0 \}
\]


\end{frame}

\begin{frame}{Examples of Sets (contd.)}
Students scoring above 80:


\[
I = \{ x \in S \mid \text{score}(x) > 80 \}
\]



Letters in the word APPLE:


\[
J = \{ x \in \{A,B,C,\dots,Z\} \mid x \text{ appears in ``APPLE''} \}
\]


\end{frame}

% Inclusion-Exclusion Principle
\begin{frame}{Inclusion-Exclusion Principle}
General Formula:


\[
\left| \bigcup_{i=1}^{\infty} A_i \right|
= \sum_{k=1}^{\infty} (-1)^{k+1} 
\left( \sum_{1 \leq i_1 < i_2 < \cdots < i_k} 
\big| A_{i_1} \cap A_{i_2} \cap \cdots \cap A_{i_k} \big| \right)
\]


\end{frame}

\begin{frame}{Proof Sketch}
1. Finite Truncation:


\[
\left| \bigcup_{i=1}^{n} A_i \right|
= \sum_{k=1}^{n} (-1)^{k+1} 
\sum_{1 \leq i_1 < \cdots < i_k \leq n} 
\big| A_{i_1} \cap \cdots \cap A_{i_k} \big|
\]



2. Limit Process: As $n \to \infty$, union approaches infinite union.
\end{frame}

\begin{frame}{Proof Sketch (contd.)}
3. Element-wise Counting:


\[
\binom{m}{1} - \binom{m}{2} + \binom{m}{3} - \cdots + (-1)^{m+1}\binom{m}{m} = 1
\]



4. Conclusion: Each element counted exactly once.
\end{frame}

% Types of Sets
\begin{frame}{Types of Sets}
\begin{tabular}{|>{\bfseries}m{4cm}|m{6cm}|}
\hline
Type of Set & Example \\
\hline
Numerical & Sensor readings: $\{23.1, 24.5, 25.0\}$ \\
Categorical & Sentiment labels: $\{\text{Positive}, \text{Negative}\}$ \\
Time Series & Stock prices: $\{(t_1,100), (t_2,102)\}$ \\
Ordinal & Satisfaction: $\{\text{Poor}, \text{Fair}, \text{Good}\}$ \\
\hline
\end{tabular}
\end{frame}

\begin{frame}{Types of Sets (contd.)}
\begin{tabular}{|>{\bfseries}m{4cm}|m{6cm}|}
\hline
Type of Set & Example \\
\hline
Image Set & MNIST digits $\{0,1,2,\dots,9\}$ \\
Text Set & NLP corpus: $\{\text{``AI is powerful''}\}$ \\
Web/API Set & Tweets: $\{\text{Tweet}_1, \text{Tweet}_2\}$ \\
Multivariate & Patient records: $\{(\text{Age}, \text{BP})\}$ \\
\hline
\end{tabular}
\end{frame}

% Posets
\begin{frame}{Partially Ordered Sets (Posets)}
Definition: A poset $(P, \leq)$ satisfies:
\begin{itemize}
    \item Reflexive: $a \leq a$
    \item Antisymmetric: $a \leq b$ and $b \leq a \Rightarrow a=b$
    \item Transitive: $a \leq b, b \leq c \Rightarrow a \leq c$
\end{itemize}
Not all elements need to be comparable.
\end{frame}

\begin{frame}{Poset Example: Divisibility on $\{1,2,4\}$}


\[
1 \leq 2 \leq 4
\]


\begin{tikzpicture}[scale=1]
\node (1) {1};
\node (2) [above=of 1] {2};
\node (4) [above=of 2] {4};
\draw (1) -- (2);
\draw (2) -- (4);
\end{tikzpicture}
\end{frame}

\begin{frame}{Poset Example: Power set of $\{a,b\}$}


\[
\emptyset \subseteq \{a\}, \{b\} \subseteq \{a,b\}
\]


\begin{tikzpicture}[scale=1]
\node (empty) {$\emptyset$};
\node (a) [above left=of empty] {$\{a\}$};
\node (b) [above right=of empty] {$\{b\}$};
\node (ab) [above=of empty,yshift=2cm] {$\{a,b\}$};
\draw (empty) -- (a);
\draw (empty) -- (b);
\draw (a) -- (ab);
\draw (b) -- (ab);
\end{tikzpicture}
\end{frame}

\begin{frame}{Poset Example: Divisibility on $\{1,2,3,6\}$}


\[
1 \leq 2 \leq 6, \quad 1 \leq 3 \leq 6
\]


\begin{tikzpicture}[scale=1]
\node (1) {1};
\node (2) [above left=of 1] {2};
\node (3) [above right=of 1] {3};
\node (6) [above=of 1,yshift=2cm] {6};
\draw (1) -- (2);
\draw (1) -- (3);
\draw (2) -- (6);
\draw (3) -- (6);
\end{tikzpicture}
\end{frame}

\end{document}
